% preamble
\documentclass[11pt,a4paper,twocolumn]{article}
\usepackage[norsk]{babel}
\usepackage[a4paper,
            top=2cm,
            bottom=2cm,
            left=3cm,
            right=3cm,
            marginparwidth=1.75cm]{geometry}
\usepackage{amsmath}
\usepackage[norsk]{cleveref}
\usepackage[shortlabels]{enumitem}
\usepackage{graphicx}
\usepackage[T1]{fontenc}
\usepackage[hyphens]{url}
\usepackage{lipsum} 
\usepackage{siunitx}
\sisetup{
  output-complex-root = \ensuremath{\mathit{j}},
  complex-root-position = before-number,
  exponent-product = \cdot,
  output-decimal-marker = {,}
}
\usepackage{tikz}
\usepackage{circuitikz}
\usetikzlibrary{arrows.meta,positioning,calc}
\usepackage{pgfplots}
\pgfplotsset{compat=1.18}
\setlength{\parindent}{0pt}

% metadata
\title{Rapport - Obligatorisk Øving 3 (ELI1300)}
\author{
	Gruppe ADK57
	\\\\
	Odin Heggvold Bekkelund, Caritha Lorenza Furnes,
	\\
	Geir Ove Venstad og Markus Øfjord
}
\date{Høst 2026}

% definer en url til senere?
\urldef{\myoverleafurl}\url{https://tinyurl.com/2yfbtdez}

% selve dokumentet starter
\begin{document}
\maketitle % Setter sammen \title og \author og \date

\begin{abstract}
Sammendrag
\end{abstract}

% Innholdsfortegnelse automatisk generert:
%\tableofcontents

\section{Innledning}
Innledning

\subsection{Utstyrsliste}
innhold

\section{Teoretisk del}
innledning

$x$, $y$ og $z$ brukes som innganger og $F$ som utgang for alle de forskjellige portene.

\subsection{NAND}
innhold

\begin{table}[h]
	\centering % sentrer
	\setlength{\belowcaptionskip}{6pt} % margin mellom caption og tabellen
	\setlength{\tabcolsep}{3pt} % gjør tabellen mindre bred hvis det trengs

	\caption{Sannhetstabell for NAND}
	\label{tab:truth-NAND}

	\noindent
	\begin{tabular}{c|c|c} % center center center
		\hline
		$x$ & $y$ & $F$ \\
		\hline
		0 & 0 & 1 \\
		0 & 1 & 1 \\
		1 & 0 & 1 \\
		1 & 1 & 0 \\
		\hline
	\end{tabular}
\end{table}

\subsection{NOT}
innhold

\begin{table}[h]
	\centering % sentrer
	\setlength{\belowcaptionskip}{6pt} % margin mellom caption og tabellen
	\setlength{\tabcolsep}{3pt} % gjør tabellen mindre bred hvis det trengs

	\caption{Sannhetstabell for NOT}
	\label{tab:truth-NOT}

	\noindent
	\begin{tabular}{c|c} % center center center
		\hline
		$x$ & $F$ \\
		\hline
		0 & 1 \\
		1 & 0 \\
		\hline
	\end{tabular}
\end{table}

\subsection{AND}
innhold

\begin{table}[h]
	\centering % sentrer
	\setlength{\belowcaptionskip}{6pt} % margin mellom caption og tabellen
	\setlength{\tabcolsep}{3pt} % gjør tabellen mindre bred hvis det trengs

	\caption{Sannhetstabell for AND}
	\label{tab:truth-AND}

	\noindent
	\begin{tabular}{c|c|c} % center center center
		\hline
		$x$ & $y$ & $F$ \\
		\hline
		0 & 0 & 0 \\
		0 & 1 & 0 \\
		1 & 0 & 0 \\
		1 & 1 & 1 \\
		\hline
	\end{tabular}
\end{table}

\subsection{OR}
innhold

\begin{table}[h]
	\centering % sentrer
	\setlength{\belowcaptionskip}{6pt} % margin mellom caption og tabellen
	\setlength{\tabcolsep}{3pt} % gjør tabellen mindre bred hvis det trengs

	\caption{Sannhetstabell for OR}
	\label{tab:truth-OR}

	\noindent
	\begin{tabular}{c|c|c} % center center center
		\hline
		$x$ & $y$ & $F$ \\
		\hline
		0 & 0 & 0 \\
		0 & 1 & 1 \\
		1 & 0 & 1 \\
		1 & 1 & 1 \\
		\hline
	\end{tabular}
\end{table}

\subsection{XOR}
innhold

\begin{table}[h]
	\centering % sentrer
	\setlength{\belowcaptionskip}{6pt} % margin mellom caption og tabellen
	\setlength{\tabcolsep}{3pt} % gjør tabellen mindre bred hvis det trengs

	\caption{Sannhetstabell for XOR}
	\label{tab:truth-XOR}

	\noindent
	\begin{tabular}{c|c|c} % center center center
		\hline
		$x$ & $y$ & $F$ \\
		\hline
		0 & 0 & 0 \\
		0 & 1 & 1 \\
		1 & 0 & 1 \\
		1 & 1 & 0 \\
		\hline
	\end{tabular}
\end{table}

\subsection{3-inngangs NAND}
innhold

\begin{table}[h]
	\centering % sentrer
	\setlength{\belowcaptionskip}{6pt} % margin mellom caption og tabellen
	\setlength{\tabcolsep}{3pt} % gjør tabellen mindre bred hvis det trengs

	\caption{Sannhetstabell for 3-inngangs NAND}
	\label{tab:truth-triple-NAND}

	\noindent
	\begin{tabular}{c|c|c|c} % center center center
		\hline
		$x$ & $y$ & $z$ & $F$ \\
		\hline
		0 & 0 & 0 & 1 \\
		0 & 0 & 1 & 1 \\
		0 & 1 & 0 & 1 \\
		0 & 1 & 1 & 1 \\
		1 & 0 & 0 & 1 \\
		1 & 0 & 1 & 1 \\
		1 & 1 & 0 & 1 \\
		1 & 1 & 1 & 0 \\
		\hline
	\end{tabular}
\end{table}

\subsection{XNOR}
innhold

\begin{table}[h]
	\centering % sentrer
	\setlength{\belowcaptionskip}{6pt} % margin mellom caption og tabellen
	\setlength{\tabcolsep}{3pt} % gjør tabellen mindre bred hvis det trengs

	\caption{Sannhetstabell for XNOR}
	\label{tab:truth-XNOR}

	\noindent
	\begin{tabular}{c|c|c} % center center center
		\hline
		$x$ & $y$ & $F$ \\
		\hline
		0 & 0 & 1 \\
		0 & 1 & 0 \\
		1 & 0 & 0 \\
		1 & 1 & 1 \\
		\hline
	\end{tabular}
\end{table}

\section{Praktisk arbeid}
innledning

\subsection{Første ting}
innhold

\subsection{Andre ting}
innhold

\subsection{Tredje ting}
innhold

\section{Resultater \& Diskusjon}
innhold

\subsection{Realisering av XNOR ved hjelp av 74F00}
XNOR-porten kan ikke realiseres med en enkelt 74F00-IC, fordi XNOR-porten krever 5 NAND-porter for å realisere. 74F00 inneholder kun 4 NAND-porter. Den kan realiseres med to 74F00-ICer, eller ved hjelp av en 74F00 og en lignende IC med minst 2 NOT-porter. Dette er fordi 2 av NAND-portene kan erstattes med NOT-porter. (kanskje vi skal ha tegning av dette også?)

\section{Konklusjon}
innhold

referansebruk: noe står i læreboka \cite{laerebok}

\begin{thebibliography}{9}
\bibitem{laerebok}
Nilsson og Riedel \textit{Electric Circuits}, 11. Utgave, Pearson.
\bibitem{analogtech}
\url{https://www.analogtechnologies.com/document/white_paper/AWPLED-04-LED-Forward-Voltage.pdf}
\end{thebibliography}

\end{document}
